\subsection{Aplicación de la matriz al caso base}
\label{sec:matriz-caso-base}
La matriz definida en el capítulo 3 se aplica al modelo de la corrida 09 y a las evidencias disponibles del proyecto. El resultado distingue la línea base cuantificada de los criterios externos que todavía no pueden verificarse. Esta evaluación documental se refiere al diseño modelado; no certifica el edificio ni supone comprobación en obra.

La revisión identifica una brecha verificable en E05: los 19 objetos de iluminación del IDF emplean 10, 15 o 25 W/m\textsuperscript{2}, por encima del límite de referencia de 8 W/m\textsuperscript{2} para oficinas Minergie. Las seis oficinas tienen 25 W/m\textsuperscript{2}. Esto acredita una diferencia en la potencia configurada; debe contrastarse con el proyecto eléctrico antes de corregir el CB. Las fichas LED y los controles también deben comprobarse \parencite{minergieLatamReglamento2025}.

La ausencia de un modelo EDGE impide determinar su ahorro energético, hídrico o material. Asimismo, ni las 1,68 renovaciones/hora agregadas ni la media operativa de 26,53 °C permiten resolver los criterios de aire y confort. La matriz presenta estas situaciones como evidencia parcial o insuficiente, no como cumplimiento implícito.
\clearpage
\begin{table}[H]\centering\caption{Aplicación de la matriz al CB: energía y envolvente}
{\fontsize{10}{12.5}\selectfont\singlespacing\renewcommand{\arraystretch}{1.3}\setlength{\tabcolsep}{4pt}\begin{tabular}{>{\raggedright\arraybackslash}p{0.8cm}>{\raggedright\arraybackslash}p{5.5cm}>{\raggedright\arraybackslash}p{1.4cm}>{\raggedright\arraybackslash}p{6.2cm}}\toprule ID & Resultado o evidencia & Estado & Brecha y siguiente verificación \\ \midrule
E01 & 84.759,41 kWh térmicos/año; corrida 09 y series anuales conciliadas. & LB & Priorizar reducción de cargas y contrastar con confort; el CB define el origen de la comparación. \\[5pt]
E02 & Electricidad estimada: enfriamiento 45.258,71; equipos 39.307,57; iluminación 13.895,39 kWh/año. Sin modelo EDGE. & NE & Completar usos finales, áreas y sistemas y construir referencia EDGE; no usar el ahorro térmico como ahorro EDGE. \\[5pt]
E03 & Muro principal: U unidimensional estimada 2,69 W/m\textsuperscript{2}K; referencia adoptada 0,75. & NC & Incorporar aislamiento; verificar puentes térmicos, clasificación climática y conjunto constructivo. \\[5pt]
E04 & Ganancia solar exterior: 33.884,25 kWh/año. Hay protecciones modeladas, sin cobertura eficaz calculada. & EP & Calcular por fachada y horas críticas; repetir con aleros C1/C3 y verificar luz natural. \\[5pt]
\bottomrule\end{tabular}}\fuente{Elaboración propia: IDF corrida 09, exportación diaria revisada y cálculo estático de iluminación. LB: línea base; NC: no alcanzado; EP: evidencia parcial; NE: no evaluable.}\end{table}\clearpage
\begin{table}[H]\centering\caption{Aplicación de la matriz al CB: eficiencia y ambiente interior}
{\fontsize{10}{12.5}\selectfont\singlespacing\renewcommand{\arraystretch}{1.3}\setlength{\tabcolsep}{4pt}\begin{tabular}{>{\raggedright\arraybackslash}p{0.8cm}>{\raggedright\arraybackslash}p{5.5cm}>{\raggedright\arraybackslash}p{1.4cm}>{\raggedright\arraybackslash}p{6.2cm}}\toprule ID & Resultado o evidencia & Estado & Brecha y siguiente verificación \\ \midrule
E05 & 19 zonas con objetos Lights: 10, 15 o 25 W/m\textsuperscript{2}. Todas superan 8 W/m\textsuperscript{2}. & NC & Contrastar con especificación eléctrica real; ensayar LED y controles en una nueva revisión, conservando la línea base. \\[5pt]
C01 & Temperatura operativa media 26,53 °C; 80 días con media > 28 °C. Sin evaluación ocupada completa. & EP & Calcular por zona con hipótesis declaradas; una media diaria no evalúa T01. La simulación no sustituye ensayo en obra. \\[5pt]
C02 & 1,68 renovaciones/hora, media conjunta de ventilación e infiltración; sin verificación de T4. & EP & Documentar caudal por recinto y tratamiento del aire; comprobar recuperación y ruido según guía. \\[5pt]
C03 & 68,37\% del área calculada supera FLD 2\% bajo cielo cubierto. Es una métrica estática distinta. & EP & Generar simulación anual y delimitar superficie regularmente ocupada; no convertir FLD en sDA. \\[5pt]
\bottomrule\end{tabular}}\fuente{Elaboración propia: IDF corrida 09, exportación diaria revisada y cálculo estático de iluminación. LB: línea base; NC: no alcanzado; EP: evidencia parcial; NE: no evaluable.}\end{table}\clearpage
\begin{table}[H]\centering\caption{Aplicación de la matriz al CB: agua, materiales y generación}
{\fontsize{10}{12.5}\selectfont\singlespacing\renewcommand{\arraystretch}{1.3}\setlength{\tabcolsep}{4pt}\begin{tabular}{>{\raggedright\arraybackslash}p{0.8cm}>{\raggedright\arraybackslash}p{5.5cm}>{\raggedright\arraybackslash}p{1.4cm}>{\raggedright\arraybackslash}p{6.2cm}}\toprule ID & Resultado o evidencia & Estado & Brecha y siguiente verificación \\ \midrule
A01 & No hay inventario completo de aparatos ni balance de agua del proyecto. & NE & Levantar caudales y descargas, usuarios y riego; calcular referencia y propuesta. \\[5pt]
M01 & Capas y propiedades térmicas disponibles; sin inventario ambiental verificable. & NE & Obtener cantidades y EPD o factores trazables; una cubierta reflectante no acredita bajo carbono. \\[5pt]
M02 & No se dispone de fichas de emisiones de los productos especificados. & NE & Identificar productos y ensayos; separar contenido de COV de emisiones al aire interior. \\[5pt]
E06 & No hay objetos fotovoltaicos en el IDF auditado; no se acredita instalación en el proyecto. & NE & Verificar planos eléctricos y SRE; dimensionar y documentar sistema, conexión y mantenimiento. \\[5pt]
\bottomrule\end{tabular}}\fuente{Elaboración propia: IDF corrida 09, exportación diaria revisada y cálculo estático de iluminación. LB: línea base; NC: no alcanzado; EP: evidencia parcial; NE: no evaluable.}\end{table}\clearpage
\begin{table}[H]\centering\caption{Aplicación de la matriz al CB: sistemas, operación y adaptación insular}
{\fontsize{10}{12.5}\selectfont\singlespacing\renewcommand{\arraystretch}{1.3}\setlength{\tabcolsep}{4pt}\begin{tabular}{>{\raggedright\arraybackslash}p{0.8cm}>{\raggedright\arraybackslash}p{5.5cm}>{\raggedright\arraybackslash}p{1.4cm}>{\raggedright\arraybackslash}p{6.2cm}}\toprule ID & Resultado o evidencia & Estado & Brecha y siguiente verificación \\ \midrule
E07 & Sistema de cargas ideales y electricidad estimada; sin selección final de equipos documentada. & NE & Seleccionar equipo y potencia; verificar humedad, refrigerante y recuperación sin cambiar supuestos entre casos. \\[5pt]
G01 & Archivos y corridas trazables; sin plan de medición operativa del edificio. & NE & Preparar cuadro de medidores y protocolo de seguimiento y ajuste. \\[5pt]
G02 & Clima documentado; falta contraste de proveedores, transporte, durabilidad y costo. & EP & Comparar alternativas por igual servicio y vida útil; justificar selección sin inventar valores ambientales. \\[5pt]
P01 & Patios cubiertos con rejilla y paso abierto; falta cómputo de áreas netas y funcionamiento por recinto. & EP & Medir aberturas, pérdidas por rejillas y horario; el 20\% de ventana de C1/C3 no prueba este requisito. \\[5pt]
\bottomrule\end{tabular}}\fuente{Elaboración propia: IDF corrida 09, exportación diaria revisada y cálculo estático de iluminación. LB: línea base; NC: no alcanzado; EP: evidencia parcial; NE: no evaluable.}\end{table}\clearpage
\subsection{Lectura del proyecto frente a los cinco marcos}
Minergie presenta brechas directas en la potencia de iluminación y la transmitancia del muro modelado y verificaciones pendientes de envolvente, protección solar, ventilación, refrigeración y autoproducción. El alcance mínimo seleccionado no permite emitir una conclusión sobre todos sus requisitos. Para WELL falta la evaluación de condiciones térmicas ocupadas y la evidencia de productos interiores. Para LEED, el FLD estático no resuelve el crédito de iluminación anual seleccionado. CEELA permite organizar las medidas de diseño y operación, pero todavía falta demostrar varias de ellas. EDGE requiere sus tres balances y su propio caso de referencia.

Con las reglas adoptadas se registra 1 indicador LB, 2 NC, 6 EP y 7 NE. Estos recuentos describen disponibilidad y calidad de evidencia; no equivalen a porcentajes de cumplimiento. El diagnóstico prioriza iluminación, control solar, ventilación y cargas especiales, junto con el cierre de información sobre agua, materiales y operación.
